\documentclass[11pt]{article}
\usepackage[margin=1in]{geometry}
\usepackage{amsmath,amssymb,amsthm}
\usepackage[T1]{fontenc}
\usepackage{lmodern}
\usepackage[hidelinks]{hyperref}
\newtheorem{theorem}{Theorem}
\newtheorem{lemma}{Lemma}
\newcommand{\Z}{\mathbb Z}
\newcommand{\Q}{\mathbb Q}
\title{Zero Filling and Persistent Free Subgroups\\of the Figure-Eight Knot}
\author{Alec Kriebel\\\small Independent researcher\\\small ORCID: \href{https://orcid.org/0009-0001-9320-500X}{0009-0001-9320-500X}}
\date{October 14, 2026}
\begin{document}
\maketitle
\begin{abstract}
The figure-eight knot gives a negative answer to the universal rank-two
free-subgroup question printed in Motegi's 2020 Oberwolfach contribution.
It is hyperbolic and has no finite nonmeridional surgery, but no rank-two
free subgroup of its knot group injects into every rational filling group.
Indeed, zero filling is a torus bundle, so its metabelian fundamental group
kills a nontrivial second-derived word in every such free subgroup.
We give the short obstruction and verify the no-finite-surgery hypothesis
using the classical filling classification. The ingredients are classical;
this note records their application and the question's quantifiers. It does
not resolve the search for a positive hyperbolic example. Historical priority
cannot be fully cleared, and no firstness claim is made.
\end{abstract}

\section{The question and the conclusion}
For a knot $K\subset S^3$, let $G=\pi_1(S^3\setminus\nu K)$, with meridian
$\mu$ and preferred longitude $\lambda$. For a rational slope $r=p/q$,
where $\gcd(p,q)=1$ and $q\ne0$, write
\[
 N_r=\langle\!\langle\mu^p\lambda^q\rangle\!\rangle_G
      =\ker(q_r),\qquad
 q_r:G\longrightarrow\pi_1(K(r)).
\]
Thus $0$ is included; the meridian $\infty=1/0$ is excluded.
The final question on printed p.~493 of \cite{OWR} asks whether a hyperbolic
knot with no finite nonmeridional surgery admits $H\cong F_2$ such that
\begin{equation}\label{target}
 H\cap\bigcup_{r\in\Q}N_r=\{1\}.
\end{equation}
Here ``finite surgery'' means a filling with finite fundamental group.
Condition \eqref{target} requires that \emph{each} $q_r|_H$ be injective.
Injectivity of a map into a product of filling groups would be a weaker,
different condition.

\begin{theorem}\label{main}
The figure-eight knot $4_1$ is hyperbolic and every rational filling
$4_1(r)$ has infinite fundamental group. Nevertheless, for every subgroup
$H\le\pi_1(S^3\setminus\nu4_1)$ isomorphic to $F_2$,
\[
 H\cap N_0\ne\{1\}.
\]
Consequently \eqref{target} fails for this knot.
\end{theorem}

This refutes the affirmative assertion for \emph{every} knot satisfying
the stated hypotheses. It says nothing against the existence of some other
hyperbolic knot for which \eqref{target} holds. Nor does it settle a variant
that is allowed to discard finitely many slopes.

\section{The algebraic obstruction}
Use $[u,v]=uvu^{-1}v^{-1}$ and write $J'=[J,J]$ and $J''=[J',J']$.
A group $M$ is metabelian if $M''=1$.

\begin{lemma}\label{law}
If $f:J\to M$ is a homomorphism into a metabelian group, its restriction
to every free rank-two subgroup of $J$ has nontrivial kernel.
\end{lemma}
\begin{proof}
Let $a,b$ be a free basis of such a subgroup $H$. Put
$c=[a,b]$, $d=aca^{-1}$ and $w=[c,d]\in H''$.
In the free alphabet $x=a$, $y=b$, $X=a^{-1}$, $Y=b^{-1}$, the freely
reduced form of $w$ is
\[
 \texttt{xyXYxxyXYXyxxYXX}.
\]
It is a nonempty reduced word of length $16$, so $w\ne1$ in $H$, and
therefore in $J$. But $f(w)\in M''=1$.
\end{proof}

If a closed three-manifold $T$ is a torus bundle over $S^1$, the homotopy
sequence gives
\begin{equation}\label{extension}
 1\longrightarrow\Z^2\longrightarrow\pi_1(T)
 \longrightarrow\Z\longrightarrow1.
\end{equation}
The left injection uses $\pi_2(S^1)=0$; surjectivity on the right uses the
connected fiber. Since the quotient is abelian, $\pi_1(T)'$ lies in the
abelian fiber subgroup. Hence $\pi_1(T)$ is metabelian. No monodromy matrix,
linearity theorem, or claim about a particular embedding of $H$ is needed.
Applying Lemma~\ref{law} to zero filling will prove the obstruction once
its classical topology and the surgery hypothesis are established.

\section{The required figure-eight topology}
The figure-eight exterior is hyperbolic; see the ideal-tetrahedron
construction in Thurston's notes \cite[Chapter~4]{Th}.
Brittenham--Wu \cite[Theorem~1.1(4)]{BW} record the complete rational
exceptional-slope list, attributing this case to Thurston:
\[
 E=\{0,\pm1,\pm2,\pm3,\pm4\}.
\]
Every rational slope outside $E$ gives a closed hyperbolic manifold and
therefore infinite fundamental group. Indeed, a finite fundamental group
would make its universal cover a finite-sheeted cover of a compact
manifold, whereas that cover is the noncompact space $\mathbb H^3$.

For completeness, all nine remaining slopes also have infinite groups:
\begin{center}
\begin{tabular}{c|l|l}
Slopes & Classical description & Reason for infinitude\\\hline
$0,\pm4$ & toroidal & an incompressible torus injects $\Z^2$\\
$\pm1$ & Seifert base $S^2(2,3,7)$ & hyperbolic triangle quotient\\
$\pm2$ & Seifert base $S^2(2,4,5)$ & hyperbolic triangle quotient\\
$\pm3$ & Seifert base $S^2(3,3,4)$ & hyperbolic triangle quotient
\end{tabular}
\end{center}
The toroidal classification is in \cite{BW}. The positive-slope Seifert
data are listed in \cite[Section~9.4, p.~71]{GM}. The three orbifold Euler
characteristics are, respectively,
\[
 -1+\tfrac12+\tfrac13+\tfrac17=-\tfrac1{42},\quad
 -1+\tfrac12+\tfrac14+\tfrac15=-\tfrac1{20},\quad
 -1+\tfrac13+\tfrac13+\tfrac14=-\tfrac1{12}.
\]
A Seifert group surjects onto its base orbifold group. Each of these
triangle groups is infinite, for it has the hyperbolic triangle action.
Amphichirality gives the negative-slope descriptions independently
\cite[Chapter~4]{Th}. In particular we do not use the displayed
$2$ versus $-3$ equality in the inspected version of \cite[Section~9.4]{GM}.

Finally, zero surgery is a torus bundle, as stated explicitly in Khoi's
Section~3, unnumbered figure-eight example on printed p.~524 \cite{Khoi}.
This also accords with Thurston's torus mapping-torus discussion
\cite[p.~70]{Th}; the latter's D-shaped surgery subscript is not our sole
basis for identifying slope zero. Equation~\eqref{extension} now shows
that $\pi_1(4_1(0))$ is metabelian, and Lemma~\ref{law} supplies a
nonidentity element of $H\cap\ker q_0$ for every $H\cong F_2$.
This proves Theorem~\ref{main}.
The no-finite hypothesis concerns all rational slopes; the meridional
$S^3$ filling is excluded throughout. Even imposing the absence of reducing
surgeries does not remove this example, by \cite[Lemma~2.1]{BW}.

\section{Context, credit, and limits}
The later question in \cite[Section~7.1]{IMT26} is again phrased for an
arbitrary hyperbolic knot. Its surrounding discussion motivates the search
for a positive hyperbolic example, which Theorem~\ref{main} leaves unresolved.
We do not infer an unstated intended quantifier or a zero-slope exception.

The topological classification and the group law above are classical.
Boden--Friedl \cite[Section~4.6, proof of Proposition~24, equation~(5)]{BF}
explicitly relate the genus-one free fiber, its longitude commutator, and
metabelian representations. A 2023 precursor \cite[p.~37]{IMT23} explicitly
records the figure-eight zero filling's abelian commutator subgroup.
That passage is cited as historical ingredient evidence only: the precursor's
Claim~6.3 was subsequently reported as having a gap, and no proof here relies
on that claim. The present proof instead uses the classical topology above.

Two independent primary-source priority audits found no explicit earlier
statement of this full counterexample in the text they inspected. This is
not a certificate of priority. In particular, the full body of the 2024
cyclic-persistence paper \cite{IMT24} and of the 2019 vanishing-elements paper
\cite{IMT19} was not inspected. Some longer precursor versions were scanned
and selectively read; Bart's 2001 surgery paper was compared only in relevant
publicly indexed text, not in full. Unindexed statements or remarks may
already contain this elementary obstruction. Priority cannot be fully cleared;
no claim of historical firstness, new topological ingredients, or continued
open status in 2026 is made.

The accompanying authored-only verification archive preserves the original
inputs, independent derivations and reviews, source locators and fingerprints,
and an offline replay driver for exact finite controls. These computations
supplement the proof; they neither exhaust arbitrary free subgroups nor replace
the cited topology. A pinned exact-inventory gate repairs a historical release
verifier's unanchored manifest and constant source-exclusion flag; the original
files remain unchanged. Raw third-party papers and account information are
excluded. Checksums establish consistency, not authenticity or priority.

AI tools were used extensively in solving, research, writing, reproduction,
and adversarial verification. This is an unrefereed preprint with internal AI
adversarial reviews. It has not received human peer review or proof-assistant
certification.

\begin{thebibliography}{10}\small
\bibitem{OWR} K. Motegi (with T. Ito and M. Teragaito), Dehn fillings on knot
groups, in \emph{Low-dimensional Topology}, Oberwolfach Reports 8/2020,
491--493. \href{https://doi.org/10.4171/OWR/2020/8}{doi:10.4171/OWR/2020/8}.
\bibitem{Th} W. P. Thurston, \emph{The Geometry and Topology of Three-Manifolds},
1980 lecture notes, electronic version 1.1 (2002), Chapter~4.
\url{https://library.slmath.org/nonmsri/gt3m/PDF/4.pdf}.
\bibitem{BW} M. Brittenham and Y.-Q. Wu, The classification of exceptional
Dehn surgeries on 2-bridge knots, \emph{Comm. Anal. Geom.} 9 (2001), 97--113.
\href{https://doi.org/10.4310/cag.2001.v9.n1.a4}{doi:10.4310/cag.2001.v9.n1.a4}.
The inspected author preprint is headed 1995.
\bibitem{Khoi} V. T. Khoi, On the integral of $\log x\,dy/y-\log y\,dx/x$
over the $A$-polynomial curves, \emph{Acta Math. Vietnam.} 33 (2008), 519--528.
\url{https://math.ac.vn/public/uploads/files/0803519.pdf}.
\bibitem{GM} S. Gukov and C. Manolescu, A two-variable series for knot
complements, author version, Section~9.4.
\url{https://web.stanford.edu/~cm5/surgeries.pdf}.
\bibitem{IMT26} T. Ito, K. Motegi, and M. Teragaito, Dehn filling and the
knot group II: Ubiquity of persistent elements, arXiv:2604.01697v1 (2026).
\url{https://arxiv.org/abs/2604.01697v1}.
\bibitem{BF} H. U. Boden and S. Friedl, Metabelian $\mathrm{SL}(n,\mathbb C)$
representations of knot groups, III: deformations, \emph{Q. J. Math.} 65
(2014), 817--840. \href{https://doi.org/10.1093/qmath/hat047}{doi:10.1093/qmath/hat047}.
Inspected version: \url{https://arxiv.org/abs/1208.1708v1} (2012).
\bibitem{IMT23} T. Ito, K. Motegi, and M. Teragaito, Dehn filling trivialization
on a knot group: separation and realization, arXiv:2303.15738v1 (2023).
\url{https://arxiv.org/abs/2303.15738v1}. The v2 abstract reports the Claim~6.3 gap.
\bibitem{IMT24} T. Ito, K. Motegi, and M. Teragaito, Dehn filling and the knot
group III: cyclic persistent subgroups, \emph{Bol. Soc. Mat. Mex.} 30 (2024), 99.
\href{https://doi.org/10.1007/s40590-024-00674-9}{doi:10.1007/s40590-024-00674-9}.
Full body not inspected.
\bibitem{IMT19} K. Ichihara, K. Motegi, and M. Teragaito, Vanishing nontrivial
elements in a knot group by Dehn fillings, \emph{Topology Appl.} 264 (2019),
223--232. \href{https://doi.org/10.1016/j.topol.2019.06.023}{doi:10.1016/j.topol.2019.06.023}.
Full body not inspected.
\end{thebibliography}
\end{document}
